\section{Temat i cel}

Temat projektu to porównanie wydajności relacyjnych i nierelacyjnych systemów
zarządzania bazą danych na przykładzie systemu rezerwacji biletów na
wydarzenia. Badamy cztery silniki: PostgreSQL~17, MariaDB~11.8, MongoDB~8.0
oraz Valkey~8.1.

Celem jest zmierzenie, jak te silniki wykonują operacje CRUD na zbiorach
danych o rozmiarze od pół miliona do dziesięciu milionów biletów. Drugim celem
jest sprawdzenie, jak radzą sobie z problemem typowym dla tej domeny: wielu
klientów jednocześnie próbuje zarezerwować te same miejsca. Domena łączy
transakcje, blokady, dane półstrukturalne (metadane wydarzeń i biletów w~JSON)
i duży ruch równoległy, więc każdy z tych elementów wynika z~zadania, a~nie
został dopisany na potrzeby testu.

\begin{table}[H]
\centering
\caption{Badane SZBD i ich rola w projekcie.}
\begin{tabular}{llp{7.4cm}}
\toprule
SZBD & model danych & rola \\
\midrule
PostgreSQL 17 & relacyjny + JSONB & 12 tabel, transakcje, blokady wierszy przy rezerwacji miejsca \\
MariaDB 11.8 & relacyjny (InnoDB) + JSON & ten sam schemat, bezpośrednie porównanie dwóch silników relacyjnych \\
MongoDB 8.0 & dokumentowy & wydarzenie z osadzonym cennikiem, rezerwacja z osadzonymi biletami i płatnością \\
Valkey 8.1 & klucz--wartość w pamięci & zbiory wolnych miejsc, blokady miejsc z TTL, liczniki sprzedaży, rezerwacje jako JSON \\
\bottomrule
\end{tabular}
\end{table}

Hipotezy badawcze sformułowane we wstępnej koncepcji:

\begin{description}
  \item[H1] Wraz ze wzrostem liczby równoległych klientów rezerwujących te same
  miejsca przewaga wydajnościowa Valkey i MongoDB nad bazami relacyjnymi rośnie.
  \item[H2] (zapasowa) Różnice wydajności między silnikami rosną wraz ze
  wzrostem liczby rekordów.
\end{description}

W stosunku do wstępnej koncepcji zmieniliśmy technologię aplikacji testowej.
Zamiast skryptu w Pythonie powstała aplikacja w języku Go podzielona na
mikroserwisy: osobny serwis gRPC dla każdej bazy, serwis wykonujący plan
testów i klient wiersza poleceń. Każdą operację śledzimy przez
OpenTelemetry, a identyfikator śladu (\emph{trace id}) zapisujemy przy każdym
wyniku. Kod znajduje się w~pięciu publicznych repozytoriach: repozytorium spinającym \texttt{ticketbench} i~czterech repozytoriach komponentów~\cite{repo}.
